\documentclass[../main.tex]{subfiles}
\begin{document}

\chapter{プロセス間通信}
\lessoninfo{
  video={P3L3},
  textbook={OSTEP \cite{ostep} ch.~5；Silberschatz ら \cite{silberschatz2018} ch.~3；Kerrisk \cite{kerrisk2010} ch.~43--54，56},
  keywords={プロセス間通信，メッセージパッシング，ポート，パイプ，メッセージキュー，ソケット，共有メモリ，コピーとマップ，ローカル手続き呼び出し，共有メモリセグメント，System V IPC，POSIX 共有メモリ，メモリマップドファイル，プロセス共有ミューテックス},
}

第2章では，それぞれ独自のアドレス空間に隔離されたプロセスがデータをやり取り
する方法としてプロセス間通信（IPC）を導入し，メッセージパッシングと共有
メモリを区別した．この章では，オペレーティングシステムが両者のために提供する
機構を調べる．共有メモリは第8章のアドレス変換の上に成り立つ．アドレス変換に
よってオペレーティングシステムは同じ物理メモリを複数のアドレス空間に対応
付けられる．共有メモリはやり取りのたびにカーネルが介在することを避けるが，
メッセージパッシングではカーネルが担う同期と通信プロトコルをプロセスに委ねる．

\section{プロセス間通信}
\label{sec:l09-ipc}

IPC とは，プロセスどうしの相互作用のためにオペレーティングシステムが
サポートする機構の総称であり，あるアドレス空間から別のアドレス空間へデータを
移す\emph{通信}と，プロセスの動作の\emph{協調}からなる．\source{P3L3，プロセス
間通信；Silberschatz ら ch.~3．}\index{ぷろせすかんつうしん@プロセス間通信}%
これらの機構は次の4つに大別される．
\begin{description}
  \item[メッセージパッシング] プロセスは，カーネルが維持するチャネルを通じて
    メッセージをやり取りする．ソケット，パイプ，メッセージキューなどである．
  \item[メモリベースの IPC] オペレーティングシステムが複数のプロセスからメモリ
    にアクセスできるようにし，プロセスはそれを読み書きして通信する．共有メモリ
    やメモリマップドファイルなどである．
  \item[より高水準のセマンティクス] データ転送により多くの構造を与える機構．
    複数のプロセスが読み書きするファイルや，遠隔手続き呼び出し（第13章）である．
  \item[同期プリミティブ] ミューテックスやセマフォなど，プロセスが協調する
    ために用いる構成要素（第10章）である．
\end{description}
この章では，メッセージパッシングと共有メモリ，そして共有メモリが必要とする
同期を中心に扱う．

\section{メッセージベースの IPC}
\label{sec:l09-message}

メッセージベースの IPC\index{めっせーじぱっしんぐ@メッセージパッシング} では，
プロセスがメッセージを作り，オペレーティングシステムが作成して維持するチャネル
---典型的にはカーネルメモリ内の FIFO キューとして構成されたバッファ---を通じて
送受信する．\source{P3L3，メッセージベースの IPC；Silberschatz ら ch.~3．}
プロセスがこのバッファに直接アクセスすることはない．

\begin{definition}[ポート]
  \term[ぽーと]{ポート}[port]とは，プロセスがメッセージパッシングのチャネルを
  利用するためのインタフェースである．プロセスはポートへメッセージを送り
  （書き），ポートからメッセージを受け取り（読み），カーネルがポートの間で
  メッセージを移動させる．
\end{definition}

カーネルは，通信の確立にも，個々の IPC 操作のすべてにも関与する．送信は，
メッセージを送信側のアドレス空間からチャネルへコピーするシステムコールであり，
受信は，それをチャネルから受信側のアドレス空間へコピーするシステムコールで
ある．したがって，クライアントとサーバの間の要求・応答のやり取りには，ユーザ/
カーネル境界の横断が4回とデータコピーが4回必要になる
（\cref{fig:l09-message-ipc}）．

\begin{figure}[htbp]
  \centering
  % Message-based IPC: a request-response exchange through a kernel channel.
% Every numbered crossing is one system call and one data copy.
% Usage: % Message-based IPC: a request-response exchange through a kernel channel.
% Every numbered crossing is one system call and one data copy.
% Usage: % Message-based IPC: a request-response exchange through a kernel channel.
% Every numbered crossing is one system call and one data copy.
% Usage: \input{figures/l09-message-ipc}   (width ~116 mm)
\begin{tikzpicture}[x=1mm, y=1mm, font=\footnotesize\sffamily,
  proc/.style={draw, rounded corners=2pt, fill=ln-box, align=center,
               minimum width=26mm, minimum height=9mm, inner sep=1pt},
  port/.style={draw, fill=ln-accent!15, minimum width=4mm,
               minimum height=21mm, inner sep=0pt},
  cell/.style={draw, fill=white, minimum width=4mm, minimum height=4.5mm,
               inner sep=0pt},
  msgc/.style={cell, fill=ln-accent!35},
  stepn/.style={circle, draw, fill=white, inner sep=0.5pt,
               minimum size=3.6mm, font=\scriptsize\sffamily\bfseries},
  lab/.style={font=\scriptsize\sffamily, text=ln-muted},
  >={Stealth[length=1.6mm]}]
  % processes
  \node[proc] (p1) at (14,28) {P1 (\iflangja{クライアント}{client})};
  \node[proc] (p2) at (98,28) {P2 (\iflangja{サーバ}{server})};
  % user/kernel boundary
  \draw[dashed, ln-rule, thick] (0,19) -- (116,19);
  \node[lab, anchor=south] at (56,19.5) {\iflangja{ユーザモード}{user mode}};
  \node[lab, anchor=north] at (56,18.5) {\iflangja{カーネルモード}{kernel mode}};
  % channel
  \draw[rounded corners=2pt, fill=ln-box] (26,-20) rectangle (86,10);
  \node[port] (pt1) at (32,-3.5) {};
  \node[port] (pt2) at (80,-3.5) {};
  \node[lab, anchor=north] at (32,-14.6) {\iflangja{ポート}{port}};
  \node[lab, anchor=north] at (80,-14.6) {\iflangja{ポート}{port}};
  \foreach \i in {0,...,7} {
    \pgfmathsetmacro{\x}{42+4*\i}
    \ifnum\i>5 \node[msgc] at (\x,1.5) {}; \else \node[cell] at (\x,1.5) {}; \fi
    \ifnum\i<2 \node[msgc] at (\x,-8.5) {}; \else \node[cell] at (\x,-8.5) {}; \fi
  }
  \draw[->] (34,1.5) -- (40,1.5);
  \draw[->] (72,1.5) -- (78,1.5);
  \draw[->] (78,-8.5) -- (72,-8.5);
  \draw[->] (40,-8.5) -- (34,-8.5);
  \node[lab, anchor=south] at (56,4.2) {\iflangja{要求}{requests}};
  \node[lab, anchor=north] at (56,-11.2) {\iflangja{応答}{responses}};
  \node[lab, anchor=north] at (56,-20.5)
    {\iflangja{チャネル（カーネル内のバッファ）}{channel (kernel buffers)}};
  % the four system calls
  \draw[->, thick] (22,23.5) -- (22,1.5) -- (30,1.5);
  \draw[->, thick] (82,1.5) -- (90,1.5) -- (90,23.5);
  \draw[->, thick] (106,23.5) -- (106,-8.5) -- (82,-8.5);
  \draw[->, thick] (30,-8.5) -- (6,-8.5) -- (6,23.5);
  \node[stepn] at (22,19) {1};
  \node[stepn] at (90,19) {2};
  \node[stepn] at (106,19) {3};
  \node[stepn] at (6,19) {4};
  \node[lab, anchor=west, text=black] at (24.2,13) {\texttt{send}};
  \node[lab, anchor=east, text=black] at (88,13)   {\texttt{recv}};
  \node[lab, anchor=east, text=black] at (104,13)  {\texttt{send}};
  \node[lab, anchor=west, text=black] at (8.2,13)  {\texttt{recv}};
\end{tikzpicture}
   (width ~116 mm)
\begin{tikzpicture}[x=1mm, y=1mm, font=\footnotesize\sffamily,
  proc/.style={draw, rounded corners=2pt, fill=ln-box, align=center,
               minimum width=26mm, minimum height=9mm, inner sep=1pt},
  port/.style={draw, fill=ln-accent!15, minimum width=4mm,
               minimum height=21mm, inner sep=0pt},
  cell/.style={draw, fill=white, minimum width=4mm, minimum height=4.5mm,
               inner sep=0pt},
  msgc/.style={cell, fill=ln-accent!35},
  stepn/.style={circle, draw, fill=white, inner sep=0.5pt,
               minimum size=3.6mm, font=\scriptsize\sffamily\bfseries},
  lab/.style={font=\scriptsize\sffamily, text=ln-muted},
  >={Stealth[length=1.6mm]}]
  % processes
  \node[proc] (p1) at (14,28) {P1 (\iflangja{クライアント}{client})};
  \node[proc] (p2) at (98,28) {P2 (\iflangja{サーバ}{server})};
  % user/kernel boundary
  \draw[dashed, ln-rule, thick] (0,19) -- (116,19);
  \node[lab, anchor=south] at (56,19.5) {\iflangja{ユーザモード}{user mode}};
  \node[lab, anchor=north] at (56,18.5) {\iflangja{カーネルモード}{kernel mode}};
  % channel
  \draw[rounded corners=2pt, fill=ln-box] (26,-20) rectangle (86,10);
  \node[port] (pt1) at (32,-3.5) {};
  \node[port] (pt2) at (80,-3.5) {};
  \node[lab, anchor=north] at (32,-14.6) {\iflangja{ポート}{port}};
  \node[lab, anchor=north] at (80,-14.6) {\iflangja{ポート}{port}};
  \foreach \i in {0,...,7} {
    \pgfmathsetmacro{\x}{42+4*\i}
    \ifnum\i>5 \node[msgc] at (\x,1.5) {}; \else \node[cell] at (\x,1.5) {}; \fi
    \ifnum\i<2 \node[msgc] at (\x,-8.5) {}; \else \node[cell] at (\x,-8.5) {}; \fi
  }
  \draw[->] (34,1.5) -- (40,1.5);
  \draw[->] (72,1.5) -- (78,1.5);
  \draw[->] (78,-8.5) -- (72,-8.5);
  \draw[->] (40,-8.5) -- (34,-8.5);
  \node[lab, anchor=south] at (56,4.2) {\iflangja{要求}{requests}};
  \node[lab, anchor=north] at (56,-11.2) {\iflangja{応答}{responses}};
  \node[lab, anchor=north] at (56,-20.5)
    {\iflangja{チャネル（カーネル内のバッファ）}{channel (kernel buffers)}};
  % the four system calls
  \draw[->, thick] (22,23.5) -- (22,1.5) -- (30,1.5);
  \draw[->, thick] (82,1.5) -- (90,1.5) -- (90,23.5);
  \draw[->, thick] (106,23.5) -- (106,-8.5) -- (82,-8.5);
  \draw[->, thick] (30,-8.5) -- (6,-8.5) -- (6,23.5);
  \node[stepn] at (22,19) {1};
  \node[stepn] at (90,19) {2};
  \node[stepn] at (106,19) {3};
  \node[stepn] at (6,19) {4};
  \node[lab, anchor=west, text=black] at (24.2,13) {\texttt{send}};
  \node[lab, anchor=east, text=black] at (88,13)   {\texttt{recv}};
  \node[lab, anchor=east, text=black] at (104,13)  {\texttt{send}};
  \node[lab, anchor=west, text=black] at (8.2,13)  {\texttt{recv}};
\end{tikzpicture}
   (width ~116 mm)
\begin{tikzpicture}[x=1mm, y=1mm, font=\footnotesize\sffamily,
  proc/.style={draw, rounded corners=2pt, fill=ln-box, align=center,
               minimum width=26mm, minimum height=9mm, inner sep=1pt},
  port/.style={draw, fill=ln-accent!15, minimum width=4mm,
               minimum height=21mm, inner sep=0pt},
  cell/.style={draw, fill=white, minimum width=4mm, minimum height=4.5mm,
               inner sep=0pt},
  msgc/.style={cell, fill=ln-accent!35},
  stepn/.style={circle, draw, fill=white, inner sep=0.5pt,
               minimum size=3.6mm, font=\scriptsize\sffamily\bfseries},
  lab/.style={font=\scriptsize\sffamily, text=ln-muted},
  >={Stealth[length=1.6mm]}]
  % processes
  \node[proc] (p1) at (14,28) {P1 (\iflangja{クライアント}{client})};
  \node[proc] (p2) at (98,28) {P2 (\iflangja{サーバ}{server})};
  % user/kernel boundary
  \draw[dashed, ln-rule, thick] (0,19) -- (116,19);
  \node[lab, anchor=south] at (56,19.5) {\iflangja{ユーザモード}{user mode}};
  \node[lab, anchor=north] at (56,18.5) {\iflangja{カーネルモード}{kernel mode}};
  % channel
  \draw[rounded corners=2pt, fill=ln-box] (26,-20) rectangle (86,10);
  \node[port] (pt1) at (32,-3.5) {};
  \node[port] (pt2) at (80,-3.5) {};
  \node[lab, anchor=north] at (32,-14.6) {\iflangja{ポート}{port}};
  \node[lab, anchor=north] at (80,-14.6) {\iflangja{ポート}{port}};
  \foreach \i in {0,...,7} {
    \pgfmathsetmacro{\x}{42+4*\i}
    \ifnum\i>5 \node[msgc] at (\x,1.5) {}; \else \node[cell] at (\x,1.5) {}; \fi
    \ifnum\i<2 \node[msgc] at (\x,-8.5) {}; \else \node[cell] at (\x,-8.5) {}; \fi
  }
  \draw[->] (34,1.5) -- (40,1.5);
  \draw[->] (72,1.5) -- (78,1.5);
  \draw[->] (78,-8.5) -- (72,-8.5);
  \draw[->] (40,-8.5) -- (34,-8.5);
  \node[lab, anchor=south] at (56,4.2) {\iflangja{要求}{requests}};
  \node[lab, anchor=north] at (56,-11.2) {\iflangja{応答}{responses}};
  \node[lab, anchor=north] at (56,-20.5)
    {\iflangja{チャネル（カーネル内のバッファ）}{channel (kernel buffers)}};
  % the four system calls
  \draw[->, thick] (22,23.5) -- (22,1.5) -- (30,1.5);
  \draw[->, thick] (82,1.5) -- (90,1.5) -- (90,23.5);
  \draw[->, thick] (106,23.5) -- (106,-8.5) -- (82,-8.5);
  \draw[->, thick] (30,-8.5) -- (6,-8.5) -- (6,23.5);
  \node[stepn] at (22,19) {1};
  \node[stepn] at (90,19) {2};
  \node[stepn] at (106,19) {3};
  \node[stepn] at (6,19) {4};
  \node[lab, anchor=west, text=black] at (24.2,13) {\texttt{send}};
  \node[lab, anchor=east, text=black] at (88,13)   {\texttt{recv}};
  \node[lab, anchor=east, text=black] at (104,13)  {\texttt{send}};
  \node[lab, anchor=west, text=black] at (8.2,13)  {\texttt{recv}};
\end{tikzpicture}

  \caption{メッセージベースの IPC による要求・応答のやり取り．番号を付した
    ユーザ/カーネル境界の横断はそれぞれ，1回のシステムコールと，プロセスの
    アドレス空間とカーネルのチャネルとの間の1回のコピーである．}
  \label{fig:l09-message-ipc}
\end{figure}

メッセージパッシングの利点は単純さである．チャネルの管理と同期はカーネルが
行う---例えば空のチャネルに対する受信は，メッセージが届くまでブロックする．
欠点はオーバーヘッドであり，境界の横断とコピーの代価をメッセージごとに支払う．

\section{メッセージパッシングの形態}
\label{sec:l09-forms}

Unix 系システムは3つの代表的なメッセージパッシングの形態を提供する．
\source{P3L3，メッセージパッシングの形態；OSTEP ch.~5；Kerrisk ch.~44，46，
56．}
それらは何を転送するかと，カーネルがどれだけの処理を行うかが異なる．

\begin{definition}[パイプ]
  \term[ぱいぷ]{パイプ}[pipe]とは，バイトストリームを運ぶ，2つの端点の間の
  チャネルである．一方の端に書かれたバイトは，メッセージ境界を持たずに，
  同じ順序で他方の端から読まれる．
\end{definition}

POSIX の \texttt{pipe(fd)} 呼び出しはパイプを作り，各端のファイルディスクリプタ
を2要素の配列 \texttt{fd} に格納する．典型的な用途は，シェルが
\texttt{ls | wc -l} に対して行うように，あるプロセスの出力を別のプロセスの
入力につなぐことである．

\begin{definition}[メッセージキュー]
  \term[めっせーじきゅー]{メッセージキュー}[message queue]とは，プロセスの
  間で離散的なメッセージを運ぶチャネルである．送信側は正しく整形された
  メッセージを渡し，受信側は完全なメッセージを受け取る．
\end{definition}

カーネルはメッセージの始まりと終わりを知っているので，メッセージを管理できる．
例えば，送信側が付けた優先度の順に配送したり（POSIX），受信側が型で
メッセージを選べるようにしたり（System~V）できる．Unix 系システムは
メッセージキューに System~V と POSIX の2つのインタフェースを提供する．

\begin{definition}[ソケット]
  \term[そけっと]{ソケット}[socket]とは，\texttt{socket()} システムコールで
  作られる通信の端点である．この呼び出しはカーネルレベルのソケットバッファを
  確保し，TCP/IP プロトコルスタックなど，その通信に必要なカーネルレベルの
  処理をそれに結び付ける．
\end{definition}

\texttt{send()} と \texttt{recv()} は，ユーザバッファからソケットバッファへ，
またソケットバッファから外へデータを渡す．ストリームソケット（TCP）は
バイトストリームを，データグラムソケット（UDP）は離散的なメッセージを配送
する．プロセスが異なるマシン上で動く場合，チャネルはプロトコルスタックを介して
プロセスをネットワークデバイスに接続する．同じマシン上で動く場合，カーネルは
完全なプロトコルスタックを迂回できる．\cref{tab:l09-forms} に3つの形態を
比較する．

\begin{table}[htbp]
  \centering
  \caption{メッセージパッシングの3つの形態．}
  \label{tab:l09-forms}
  \small
  \begin{tabular}{@{}llll@{}}
    \toprule
    & パイプ & メッセージキュー & ソケット \\
    \midrule
    転送の単位 & バイトストリーム & メッセージ & ストリーム/データグラム \\
    カーネルの処理 & バッファリング & 境界，優先度 & プロトコルスタック \\
    マシン & 同一 & 同一 & 同一または別 \\
    インタフェース & \texttt{pipe()} & System~V，POSIX & \texttt{socket()} \\
    \bottomrule
  \end{tabular}
\end{table}

\section{共有メモリ IPC}
\label{sec:l09-shm}

共有メモリ IPC では，複数のプロセスが，そのすべてからアクセスできるメモリ
領域を読み書きする．\source{P3L3，共有メモリ IPC；Silberschatz ら
ch.~3．}\index{きょうゆうめもり@共有メモリ}%
オペレーティングシステムは，同じ物理ページをプロセスの仮想アドレス空間に
対応付けることでこのチャネルを確立する．すなわち，参加するすべてのプロセスの
ページテーブルエントリが同じページフレームを指す（\cref{fig:l09-shm-mapping}）．
ページングによる対応付け一般と同様に（第8章），プロセスが領域を見る仮想
アドレスは同じである必要はなく，物理フレームも連続である必要はない．

\begin{figure}[htbp]
  \centering
  % Shared-memory IPC: the same (non-contiguous) physical frames are mapped
% into two virtual address spaces at different virtual addresses.
% Usage: % Shared-memory IPC: the same (non-contiguous) physical frames are mapped
% into two virtual address spaces at different virtual addresses.
% Usage: % Shared-memory IPC: the same (non-contiguous) physical frames are mapped
% into two virtual address spaces at different virtual addresses.
% Usage: \input{figures/l09-shm-mapping}   (width ~116 mm)
\begin{tikzpicture}[x=1mm, y=1mm, font=\footnotesize\sffamily,
  slot/.style={draw, fill=white, minimum width=18mm, minimum height=5mm,
               inner sep=0pt, font=\scriptsize\sffamily},
  sh/.style={slot, fill=ln-accent!20},
  oth/.style={slot, fill=ln-box},
  hd/.style={font=\scriptsize\sffamily, align=center, text width=30mm},
  lab/.style={font=\scriptsize\sffamily, text=ln-muted},
  >={Stealth[length=1.6mm]}]
  \foreach \i in {0,...,7} {
    \node[slot] (a\i) at (17, 2.5+5*\i) {};
    \node[slot] (m\i) at (59, 2.5+5*\i) {};
    \node[slot] (b\i) at (101, 2.5+5*\i) {};
  }
  % other (private) frames
  \foreach \i in {0,3,4,7} { \node[oth] at (59, 2.5+5*\i) {}; }
  % shared pages
  \node[sh] (a5) at (17,27.5)  {\iflangja{共有 0}{shared 0}};
  \node[sh] (a6) at (17,32.5)  {\iflangja{共有 1}{shared 1}};
  \node[sh] (m6) at (59,32.5)  {\iflangja{共有 0}{shared 0}};
  \node[sh] (m2) at (59,12.5)  {\iflangja{共有 1}{shared 1}};
  \node[sh] (b1) at (101,7.5)  {\iflangja{共有 0}{shared 0}};
  \node[sh] (b2) at (101,12.5) {\iflangja{共有 1}{shared 1}};
  \draw[->] (a5.east) -- (m6.west);
  \draw[->] (a6.east) -- (m2.west);
  \draw[->] (b1.west) -- (m6.east);
  \draw[->] (b2.west) -- (m2.east);
  % headers
  \node[hd, anchor=south] at (17,41)  {\iflangja{P1 の仮想アドレス空間}{virtual address\\space of P1}};
  \node[hd, anchor=south] at (59,41)  {\iflangja{物理メモリ}{physical memory}};
  \node[hd, anchor=south] at (101,41) {\iflangja{P2 の仮想アドレス空間}{virtual address\\space of P2}};
  % virtual addresses of the region
  \draw[ln-muted] (6,25) -- (8,25);
  \node[lab, anchor=east] at (6,25) {$\mathit{VA}_1$};
  \draw[ln-muted] (110,5) -- (112,5);
  \node[lab, anchor=west] at (112,5) {$\mathit{VA}_2$};
\end{tikzpicture}
   (width ~116 mm)
\begin{tikzpicture}[x=1mm, y=1mm, font=\footnotesize\sffamily,
  slot/.style={draw, fill=white, minimum width=18mm, minimum height=5mm,
               inner sep=0pt, font=\scriptsize\sffamily},
  sh/.style={slot, fill=ln-accent!20},
  oth/.style={slot, fill=ln-box},
  hd/.style={font=\scriptsize\sffamily, align=center, text width=30mm},
  lab/.style={font=\scriptsize\sffamily, text=ln-muted},
  >={Stealth[length=1.6mm]}]
  \foreach \i in {0,...,7} {
    \node[slot] (a\i) at (17, 2.5+5*\i) {};
    \node[slot] (m\i) at (59, 2.5+5*\i) {};
    \node[slot] (b\i) at (101, 2.5+5*\i) {};
  }
  % other (private) frames
  \foreach \i in {0,3,4,7} { \node[oth] at (59, 2.5+5*\i) {}; }
  % shared pages
  \node[sh] (a5) at (17,27.5)  {\iflangja{共有 0}{shared 0}};
  \node[sh] (a6) at (17,32.5)  {\iflangja{共有 1}{shared 1}};
  \node[sh] (m6) at (59,32.5)  {\iflangja{共有 0}{shared 0}};
  \node[sh] (m2) at (59,12.5)  {\iflangja{共有 1}{shared 1}};
  \node[sh] (b1) at (101,7.5)  {\iflangja{共有 0}{shared 0}};
  \node[sh] (b2) at (101,12.5) {\iflangja{共有 1}{shared 1}};
  \draw[->] (a5.east) -- (m6.west);
  \draw[->] (a6.east) -- (m2.west);
  \draw[->] (b1.west) -- (m6.east);
  \draw[->] (b2.west) -- (m2.east);
  % headers
  \node[hd, anchor=south] at (17,41)  {\iflangja{P1 の仮想アドレス空間}{virtual address\\space of P1}};
  \node[hd, anchor=south] at (59,41)  {\iflangja{物理メモリ}{physical memory}};
  \node[hd, anchor=south] at (101,41) {\iflangja{P2 の仮想アドレス空間}{virtual address\\space of P2}};
  % virtual addresses of the region
  \draw[ln-muted] (6,25) -- (8,25);
  \node[lab, anchor=east] at (6,25) {$\mathit{VA}_1$};
  \draw[ln-muted] (110,5) -- (112,5);
  \node[lab, anchor=west] at (112,5) {$\mathit{VA}_2$};
\end{tikzpicture}
   (width ~116 mm)
\begin{tikzpicture}[x=1mm, y=1mm, font=\footnotesize\sffamily,
  slot/.style={draw, fill=white, minimum width=18mm, minimum height=5mm,
               inner sep=0pt, font=\scriptsize\sffamily},
  sh/.style={slot, fill=ln-accent!20},
  oth/.style={slot, fill=ln-box},
  hd/.style={font=\scriptsize\sffamily, align=center, text width=30mm},
  lab/.style={font=\scriptsize\sffamily, text=ln-muted},
  >={Stealth[length=1.6mm]}]
  \foreach \i in {0,...,7} {
    \node[slot] (a\i) at (17, 2.5+5*\i) {};
    \node[slot] (m\i) at (59, 2.5+5*\i) {};
    \node[slot] (b\i) at (101, 2.5+5*\i) {};
  }
  % other (private) frames
  \foreach \i in {0,3,4,7} { \node[oth] at (59, 2.5+5*\i) {}; }
  % shared pages
  \node[sh] (a5) at (17,27.5)  {\iflangja{共有 0}{shared 0}};
  \node[sh] (a6) at (17,32.5)  {\iflangja{共有 1}{shared 1}};
  \node[sh] (m6) at (59,32.5)  {\iflangja{共有 0}{shared 0}};
  \node[sh] (m2) at (59,12.5)  {\iflangja{共有 1}{shared 1}};
  \node[sh] (b1) at (101,7.5)  {\iflangja{共有 0}{shared 0}};
  \node[sh] (b2) at (101,12.5) {\iflangja{共有 1}{shared 1}};
  \draw[->] (a5.east) -- (m6.west);
  \draw[->] (a6.east) -- (m2.west);
  \draw[->] (b1.west) -- (m6.east);
  \draw[->] (b2.west) -- (m2.east);
  % headers
  \node[hd, anchor=south] at (17,41)  {\iflangja{P1 の仮想アドレス空間}{virtual address\\space of P1}};
  \node[hd, anchor=south] at (59,41)  {\iflangja{物理メモリ}{physical memory}};
  \node[hd, anchor=south] at (101,41) {\iflangja{P2 の仮想アドレス空間}{virtual address\\space of P2}};
  % virtual addresses of the region
  \draw[ln-muted] (6,25) -- (8,25);
  \node[lab, anchor=east] at (6,25) {$\mathit{VA}_1$};
  \draw[ln-muted] (110,5) -- (112,5);
  \node[lab, anchor=west] at (112,5) {$\mathit{VA}_2$};
\end{tikzpicture}

  \caption{共有メモリ．2つのページフレームが両方のアドレス空間に対応付け
    られている．P1 はこの領域を仮想アドレス $\mathit{VA}_1$ に，P2 は別の
    アドレス $\mathit{VA}_2$ に見ており，フレームは物理メモリ上で隣接して
    いない．}
  \label{fig:l09-shm-mapping}
\end{figure}

\caution{$\mathit{VA}_1 \neq \mathit{VA}_2$ であるから，一方のプロセスが
領域に格納したポインタは他方では有効でない．}

対応付けができてしまえば，オペレーティングシステムは介在しない．読み書きは
通常のメモリアクセスであり，システムコールが必要なのは設定のときだけである．
データコピーは削減されうるが，なくなるわけではない．プロセスが共有領域の中で
生成したデータは直ちに他のプロセスから見えるが，そのアドレス空間の別の場所に
すでに存在するデータは，やはり領域へコピーしなければならない．

代償として，オペレーティングシステムはメッセージパッシングに対して提供して
いたサービスをもはや行わない．プロセスは並行アクセスを明示的に同期しなければ
ならず（\cref{sec:l09-sync}），通信プロトコル---データの形式，受信側が
データの到着をどう知るか，領域内の空間をどう管理するか---はプログラマの責任
である（\cref{sec:l09-design}）．Unix と Linux は，System~V 共有メモリ
（\cref{sec:l09-sysv}），POSIX 共有メモリ（\cref{sec:l09-posix}），および
メモリマップドファイルによって共有メモリをサポートする．Android は独自の機構
ashmem を加えている．

\section{コピーとマップ}
\label{sec:l09-copy-map}

メッセージパッシングと共有メモリは，あるアドレス空間から別のアドレス空間へ
データを移すという同じ目的を追うが，CPU 時間を費やす操作が異なる．
\source{P3L3，コピーとマップ；Silberschatz ら ch.~3．}
\begin{description}
  \item[コピー（メッセージパッシング）] CPU 時間は，データをポートへ，そして
    ポートから外へコピーすることに費やされ，送信と受信のたびに---すなわち
    メッセージごとに---システムコールを伴う．
  \item[マップ（共有メモリ）] CPU 時間は，物理メモリをアドレス空間に対応付ける
    こと---領域を設定するときの1回だけ---と，領域の中で生成されなかったデータ
    を領域へコピーすることに費やされる．
\end{description}

対応付けは比較的高価である．1回設定して多数の転送に使えば，第2章の例が示した
ように元が取れる．大きなデータであれば，1回の転送でも元が取れうる．大量の
データをあるアドレス空間から別のアドレス空間へコピーする時間は，それを対応
付ける時間よりはるかに長いからである: $t_{\mathrm{copy}} \gg
t_{\mathrm{map}}$．単純なコストモデルによって，この2つの主張を精密にできる．
システムコールのコストを $t_{\mathrm{sys}}$，コピーのコストを1バイト当たり $c$
とし，
$S$ バイトの転送があるプロセスから別のプロセスへ $n$ 回行われるとする．
メッセージパッシングは転送ごとに2回のシステムコールと2回のコピーを必要と
する．共有メモリは1回限りの設定コスト $t_{\mathrm{setup}}$ を必要とし，転送
ごとには，領域への1回のコピー---受信側はそこからデータをその場で読む---と，
プロセスがアクセスの同期と互いへの通知に費やす時間 $t_{\mathrm{sync}}$
（\cref{sec:l09-sync}）を必要とする．
\begin{equation}
  T_{\mathrm{copy}}(n) = n\,(2\,t_{\mathrm{sys}} + 2cS) ,
  \qquad
  T_{\mathrm{map}}(n) = t_{\mathrm{setup}} + n\,(cS + t_{\mathrm{sync}}) .
  \label{eq:l09-copy-map}
\end{equation}
$t_{\mathrm{sync}} < 2\,t_{\mathrm{sys}} + cS$ であれば，共有メモリは
$n > n^{*} = t_{\mathrm{setup}} / (2\,t_{\mathrm{sys}} + cS -
t_{\mathrm{sync}})$ のとき安く，そうでなければ決して追いつかない．競合の
ないプロセス共有ミューテックスが領域を保護し，受信側がブロックしない場合，
同期コストはほぼゼロであるが，ブロックして起こされるプロセスはシステム
コールを追加する．設定コストは，固定部分と，対応付けるページ数に比例する
部分からなる．1ページを対応付けるコストがその内容をコピーするコストより
小さい限り，$n^{*}$ は $S$ が大きくなるにつれて下がり，大きなデータでは1を
下回る．

\begin{example}
  \label{ex:l09-copy-map}
  $t_{\mathrm{sys}} = \SI{1}{\micro\second}$，$c$ を1キロバイト当たり
  \SI{0.1}{\micro\second} とし，設定コストを \SI{50}{\micro\second} に
  \SI{4}{\kilo\byte} ページ当たり \SI{0.1}{\micro\second} を加えたものと
  する．$S = \SI{40}{\mega\byte}$（$10^4$ ページ）の1回の転送は，メッセージ
  パッシングでは $2 + 2 \times 4000 = \SI{8002}{\micro\second}$ かかるが，
  $t_{\mathrm{sync}} \approx 0$ とすると共有メモリでは $1050 + 4000 =
  \SI{5050}{\micro\second}$ であり，送信側がデータを領域の中で直接生成する
  なら設定の \SI{1050}{\micro\second} だけで済む．$n^{*} = 1050/4002 \approx
  0.26$ である．ブロックする待機と起床（$t_{\mathrm{sync}} =
  \SI{2}{\micro\second}$）があってもこれは変わらない（$n^{*} = 1050/4000
  \approx 0.26$）が，$S = \SI{4}{\kilo\byte}$ では $n^{*}$ が $50.1/2.4
  \approx 21$ から $50.1/0.4 \approx 125$ へ上がる．
\end{example}

Windows は，同じマシン上のプロセスの間の通信のための機構である
\term[ろーかるてつづきよびだし]{ローカル手続き呼び出し}[local procedure call]
（LPC）\index{LPC|see{ローカル手続き呼び出し}}でこのトレードオフを利用して
いる．小さなメッセージはポートを介してコピーし，より大きなデータに対しては
共有メモリ領域を両方のプロセスに対応付け，それを通じてデータをやり取りする．

\begin{keypoint}
  メッセージパッシングは単純である---カーネルがチャネルを管理し，プロセスを
  同期する---が，メッセージごとにシステムコールと2回のコピーの代価を払う．
  共有メモリは1回限りの対応付けのコストを払った後はカーネルを介さずに済み，
  その代わりに明示的な同期とプログラマが定めるプロトコルを要する．小さな
  メッセージはコピーし，大きなデータや多数の転送に再利用される領域は対応
  付けるとよい．
\end{keypoint}

\section{System V 共有メモリ}
\label{sec:l09-sysv}

最も古い Unix の IPC インタフェースの1つである System~V の
インタフェース\source{P3L3，SysV 共有メモリ，SysV 共有メモリ API；Kerrisk
ch.~45，48．}は，共有メモリをセグメントとして構成する．

\begin{definition}[共有メモリセグメント]
  \term[きょうゆうめもりせぐめんと]{共有メモリセグメント}[shared memory
  segment]とは，必ずしも連続したページからなるとは限らない物理メモリの領域で
  あって，オペレーティングシステムがシステム全体のオブジェクトとして管理し，
  任意個数のプロセスが自分のアドレス空間に対応付けられるものである．
\end{definition}

セグメントはシステム全体のオブジェクトなので，オペレーティングシステムはその
個数と合計サイズの両方を制限する．\margin{Linux ではカーネルパラメータ
\texttt{shmmni}，\texttt{shmmax}，\texttt{shmall} が上限を定め，既定では最大
4096 個のセグメントが存在できる．}セグメントは\emph{キー}で識別され，同じ
キーを使うプロセスはすべて同じセグメントにアクセスする．そのライフサイクルは
4つの操作からなる（\cref{fig:l09-sysv-lifecycle}）．
\begin{description}
  \item[作成] オペレーティングシステムがセグメントのメモリを確保し，それを
    キーに結び付ける．
  \item[接続] オペレーティングシステムがセグメントを呼び出し元プロセスの
    アドレス空間に対応付ける．すなわち，プロセスの仮想アドレスからセグメントの
    物理メモリへの有効な対応付けを確立する．複数のプロセスが同じセグメントを，
    それぞれ自分の仮想アドレスで接続できる．
  \item[切り離し] オペレーティングシステムがこれらの対応付けを無効にする．
    セグメントの仮想アドレス範囲に対するページテーブルエントリはもはや有効で
    ない．セグメント自体は残る．
  \item[破棄] セグメントを取り除くのは明示的な要求だけである．それまでは，
    どのプロセスも接続していなくても存続し，何度でも接続と切り離しができる．
\end{description}

\begin{figure}[htbp]
  \centering
  % Life cycle of a System V shared memory segment.
% Usage: % Life cycle of a System V shared memory segment.
% Usage: % Life cycle of a System V shared memory segment.
% Usage: \input{figures/l09-sysv-lifecycle}   (width ~120 mm)
\begin{tikzpicture}[x=1mm, y=1mm, font=\footnotesize\sffamily,
  st/.style={draw, rounded corners=3pt, fill=ln-box, align=center,
             minimum height=11mm, inner sep=2pt},
  api/.style={font=\scriptsize\ttfamily, align=center, inner sep=1pt},
  >={Stealth[length=1.8mm]}]
  \node[st, text width=17mm] (none) at (0,0)
    {\iflangja{セグメント\\なし}{no\\segment}};
  \node[st, text width=24mm] (free) at (48,0)
    {\iflangja{存在する\\（接続なし）}{exists,\\not attached}};
  \node[st, text width=24mm] (att) at (96,0)
    {\iflangja{$n\ge 1$ 個のプロセス\\が接続}{attached by\\$n\ge 1$ processes}};
  \node[st, text width=24mm] (mark) at (96,-30)
    {\iflangja{破棄予約済み}{marked for\\destruction}};
  \draw[->] (none) -- node[api, above] {shmget()}
                      node[api, below] {IPC\_CREAT} (free);
  \draw[->] ([yshift=2.5mm]free.east) -- node[api, above] {shmat()}
                                         ([yshift=2.5mm]att.west);
  \draw[->] ([yshift=-2.5mm]att.west) -- node[api, below] {shmdt()\\(\iflangja{最後}{last})}
                                         ([yshift=-2.5mm]free.east);
  \draw[->] ([xshift=-4mm]att.north) .. controls +(-3,12) and +(3,12) ..
    node[api, above] {shmat() / shmdt()} ([xshift=4mm]att.north);
  \draw[->] (free.north) .. controls +(0,12) and +(0,12) ..
    node[api, above] {shmctl(IPC\_RMID)} (none.north);
  \draw[->] (att) -- node[api, left=1mm] {shmctl(IPC\_RMID)} (mark);
  \draw[->] (mark.west) -| node[api, above, pos=0.3] {shmdt()\,(\iflangja{最後}{last})} (none.south);
\end{tikzpicture}
   (width ~120 mm)
\begin{tikzpicture}[x=1mm, y=1mm, font=\footnotesize\sffamily,
  st/.style={draw, rounded corners=3pt, fill=ln-box, align=center,
             minimum height=11mm, inner sep=2pt},
  api/.style={font=\scriptsize\ttfamily, align=center, inner sep=1pt},
  >={Stealth[length=1.8mm]}]
  \node[st, text width=17mm] (none) at (0,0)
    {\iflangja{セグメント\\なし}{no\\segment}};
  \node[st, text width=24mm] (free) at (48,0)
    {\iflangja{存在する\\（接続なし）}{exists,\\not attached}};
  \node[st, text width=24mm] (att) at (96,0)
    {\iflangja{$n\ge 1$ 個のプロセス\\が接続}{attached by\\$n\ge 1$ processes}};
  \node[st, text width=24mm] (mark) at (96,-30)
    {\iflangja{破棄予約済み}{marked for\\destruction}};
  \draw[->] (none) -- node[api, above] {shmget()}
                      node[api, below] {IPC\_CREAT} (free);
  \draw[->] ([yshift=2.5mm]free.east) -- node[api, above] {shmat()}
                                         ([yshift=2.5mm]att.west);
  \draw[->] ([yshift=-2.5mm]att.west) -- node[api, below] {shmdt()\\(\iflangja{最後}{last})}
                                         ([yshift=-2.5mm]free.east);
  \draw[->] ([xshift=-4mm]att.north) .. controls +(-3,12) and +(3,12) ..
    node[api, above] {shmat() / shmdt()} ([xshift=4mm]att.north);
  \draw[->] (free.north) .. controls +(0,12) and +(0,12) ..
    node[api, above] {shmctl(IPC\_RMID)} (none.north);
  \draw[->] (att) -- node[api, left=1mm] {shmctl(IPC\_RMID)} (mark);
  \draw[->] (mark.west) -| node[api, above, pos=0.3] {shmdt()\,(\iflangja{最後}{last})} (none.south);
\end{tikzpicture}
   (width ~120 mm)
\begin{tikzpicture}[x=1mm, y=1mm, font=\footnotesize\sffamily,
  st/.style={draw, rounded corners=3pt, fill=ln-box, align=center,
             minimum height=11mm, inner sep=2pt},
  api/.style={font=\scriptsize\ttfamily, align=center, inner sep=1pt},
  >={Stealth[length=1.8mm]}]
  \node[st, text width=17mm] (none) at (0,0)
    {\iflangja{セグメント\\なし}{no\\segment}};
  \node[st, text width=24mm] (free) at (48,0)
    {\iflangja{存在する\\（接続なし）}{exists,\\not attached}};
  \node[st, text width=24mm] (att) at (96,0)
    {\iflangja{$n\ge 1$ 個のプロセス\\が接続}{attached by\\$n\ge 1$ processes}};
  \node[st, text width=24mm] (mark) at (96,-30)
    {\iflangja{破棄予約済み}{marked for\\destruction}};
  \draw[->] (none) -- node[api, above] {shmget()}
                      node[api, below] {IPC\_CREAT} (free);
  \draw[->] ([yshift=2.5mm]free.east) -- node[api, above] {shmat()}
                                         ([yshift=2.5mm]att.west);
  \draw[->] ([yshift=-2.5mm]att.west) -- node[api, below] {shmdt()\\(\iflangja{最後}{last})}
                                         ([yshift=-2.5mm]free.east);
  \draw[->] ([xshift=-4mm]att.north) .. controls +(-3,12) and +(3,12) ..
    node[api, above] {shmat() / shmdt()} ([xshift=4mm]att.north);
  \draw[->] (free.north) .. controls +(0,12) and +(0,12) ..
    node[api, above] {shmctl(IPC\_RMID)} (none.north);
  \draw[->] (att) -- node[api, left=1mm] {shmctl(IPC\_RMID)} (mark);
  \draw[->] (mark.west) -| node[api, above, pos=0.3] {shmdt()\,(\iflangja{最後}{last})} (none.south);
\end{tikzpicture}

  \caption{System~V 共有メモリセグメントのライフサイクル．セグメントはそれを
    使うプロセスより長く残り，取り除かれるのは \texttt{IPC\_RMID} によってのみ
    である．接続されたまま取り除かれたセグメントは，最後のプロセスが切り離した
    時点で破棄される．}
  \label{fig:l09-sysv-lifecycle}
\end{figure}

したがってセグメントは，プロセス固有のメモリというよりファイルのように振る舞う．
終了したプロセスはセグメントから切り離されるが，セグメントは破棄されるか
システムが再起動するまで残る．\caution{どのプロセスも破棄しなかったセグメント
は，再起動までメモリを占有し続ける（\cref{sec:l09-tools}）．}

\subsection{System V の API}

このインタフェースは4つの呼び出しと，キーを生成する補助関数からなる．
\begin{description}
  \item[\texttt{shmget(key, size, flags)}] は \texttt{key} に結び付いた
    セグメントを開く．存在せず，かつ \texttt{flags} に \texttt{IPC\_CREAT} が
    含まれていれば，指定された大きさのセグメントを作成する
    （\texttt{IPC\_CREAT | IPC\_EXCL} では，セグメントがすでに存在すると
    呼び出しは失敗する）．セグメントの識別子を返す．
  \item[\texttt{ftok(pathname, proj\_id)}] は，既存のファイルの同一性
    （デバイスと i ノード番号）と，0 でない整数の下位 8 ビットからキーを
    生成する．ファイルが作り直されない限り，同じファイルと同じ番号は同じキーを
    与えるので，無関係なプロセスどうしも，ファイル名と番号を取り決めることで
    セグメントを見つけられる．
  \item[\texttt{shmat(id, addr, flags)}] はセグメントを接続し，それが対応付け
    られた仮想アドレスを返す．\texttt{addr} を \texttt{NULL} にすると
    オペレーティングシステムがアドレスを選ぶ．返される型のないポインタは，
    セグメントに置く任意のデータ構造にキャストできる．
  \item[\texttt{shmdt(addr)}] は \texttt{addr} に対応付けられたセグメントを
    切り離す．
  \item[\texttt{shmctl(id, cmd, buf)}] はセグメントに対する制御コマンドを
    オペレーティングシステムに渡す．\texttt{IPC\_RMID} はセグメントを破棄する．
\end{description}

\begin{example}
  \label{ex:l09-sysv}
  プロセス A がセグメントを作成して文字列を書き込み，後から実行されるプロセス
  B が同じキーでそのセグメントを見つけ，文字列を表示してセグメントを破棄する
  （エラー処理は省略）．
  \begin{lstlisting}[language=C, numbers=none]
/* プロセス A */
key_t key = ftok("/tmp/app", 'A');
int id = shmget(key, 4096, IPC_CREAT | 0600);
char *p = shmat(id, NULL, 0);
strcpy(p, "hello");                /* B から見える */
shmdt(p);                          /* セグメントは残る */

/* プロセス B */
key_t key = ftok("/tmp/app", 'A'); /* 同じキー */
int id = shmget(key, 4096, 0);     /* 開く */
char *q = shmat(id, NULL, 0);
printf("%s\n", q);
shmdt(q);
shmctl(id, IPC_RMID, NULL);        /* 破棄する */
\end{lstlisting}
  A の \texttt{shmdt} と B の \texttt{shmat} の間，どのプロセスもセグメントを
  接続していないが，その内容は保たれている．
\end{example}

\section{POSIX 共有メモリ}
\label{sec:l09-posix}

POSIX は第2の共有メモリインタフェースを標準化しているが，これは System~V ほど
広くサポートされてはいない．例えば Linux では，カーネル 2.4 以降で
サポートされている．\source{P3L3，POSIX 共有メモリ API；Kerrisk ch.~49，54．}
キーで識別されるセグメントの代わりに，名前で識別されファイルのようにアクセス
される共有メモリオブジェクトを用いる．Linux では，これらはメモリ上の一時
ファイルシステム（tmpfs）内のファイルである．\margin{Linux ではこの tmpfs は
\texttt{/dev/shm} にマウントされる．}プロセスは，ファイルをメモリに対応付ける
一般的な機構を通じてそれらにアクセスする．

\begin{definition}[メモリマップドファイル]
  \term[めもりまっぷどふぁいる]{メモリマップドファイル}[memory-mapped
  file]とは，その内容がプロセスの仮想アドレス空間に対応付けられ，プロセスが
  メモリの読み書きによってアクセスできるようになったファイルである．同じ
  ファイルを共有として対応付けたプロセスは，同じ物理ページを対応付け，互いの
  変更を見る．
\end{definition}

通常のファイルシステム上のメモリマップドファイルは，それ自体が IPC 機構で
あり，そのデータはファイルにも永続する．POSIX 共有メモリは，メモリ上にのみ
存在するオブジェクトに同じ機構を適用したものである．その操作は System~V の
操作に対応する（\cref{tab:l09-shm-api}）．\texttt{shm\_open(name, flags,
mode)} は共有メモリオブジェクトを作成または開き，そのファイルディスクリプタを
返す．\texttt{ftruncate} がその大きさを設定する．\texttt{MAP\_SHARED} フラグ
付きの \texttt{mmap} はオブジェクトを接続してその仮想アドレスを返し，
\texttt{munmap} はそれを切り離す．\texttt{close} は，確立済みの対応付けに
影響を与えずにファイルディスクリプタを解放し，\texttt{shm\_unlink(name)} は
名前を取り除く．どのプロセスもそれを開いておらず，対応付けてもいなくなった
時点で，オブジェクトは破棄される．

\begin{table}[htbp]
  \centering
  \caption{2つの共有メモリインタフェースの対応する操作．}
  \label{tab:l09-shm-api}
  \small
  \begin{tabular}{@{}lll@{}}
    \toprule
    操作 & System~V & POSIX \\
    \midrule
    領域の名前付け & \texttt{ftok()} によるキー & オブジェクト名（\texttt{/name}） \\
    作成・オープン & \texttt{shmget()} & \texttt{shm\_open()}，\texttt{ftruncate()} \\
    接続 & \texttt{shmat()} & \texttt{mmap()} \\
    切り離し & \texttt{shmdt()} & \texttt{munmap()} \\
    破棄 & \texttt{shmctl(IPC\_RMID)} & \texttt{shm\_unlink()} \\
    \bottomrule
  \end{tabular}
\end{table}

\section{共有メモリと同期}
\label{sec:l09-sync}

共有領域に並行してアクセスするプロセスは，1つのアドレス空間の中で状態を共有
する第3章のスレッドと同じ状況にある．同期がなければ，そのアクセスは競合状態を
なす．\source{P3L3，共有メモリと同期，IPC のための PThreads 同期，他の IPC の
ための同期；Kerrisk ch.~45--47．}利用できる機構は2種類ある．1つは，プロセスを
またいで使えるように構成された Pthreads のミューテックスや条件変数のような，
スレッドライブラリの機構である．もう1つは，メッセージキューやセマフォのような，
オペレーティングシステムが提供する IPC 機構に基づく同期である．どちらの種類
でも，プロセスは2つのことを協調させなければならない．1つは共有領域への並行
アクセスの数であり，これは相互排除が制限する．もう1つはデータが消費できるよう
になる時点であり，条件変数の場合のように，一方のプロセスがそれを待つことが
でき，他方がそれを知らせられなければならない．

\subsection{Pthreads によるプロセス間の同期}

第4章の Pthreads のミューテックスと条件変数は，2つの条件の下でプロセスを同期
できる．第1に，それらの初期化に用いる属性オブジェクト
\texttt{pthread\_mutexattr\_t} と \texttt{pthread\_condattr\_t} の
プロセス共有属性を，\texttt{PTHREAD\_PROCESS\_SHARED} に設定しなければ
ならない．既定では，作成したプロセスのスレッドだけがそれらを使える．第2に，
ミューテックスと条件変数自体が共有領域に置かれていなければならない．そう
することで，すべてのプロセスが同じデータ構造を操作する．

\needspace{9\baselineskip}
\begin{example}
  \label{ex:l09-pshared}
  あるプロセスが領域のレイアウトを定める．領域は同期構成要素から始まる:
  \begin{lstlisting}[language=C, numbers=none]
typedef struct {
  pthread_mutex_t mutex;
  pthread_cond_t  cond;
  char            data[1024];
} shm_data_t;
\end{lstlisting}
  \needspace{6\baselineskip}
  そしてセグメントを作成して接続する:
  \begin{lstlisting}[language=C, numbers=none]
key_t key = ftok("/tmp/app", 120);
int id = shmget(key, sizeof(shm_data_t),
                IPC_CREAT | IPC_EXCL | 0600);
shm_data_t *shm = shmat(id, NULL, 0);
\end{lstlisting}
  \needspace{7\baselineskip}
  さらに領域の中でプロセス共有のミューテックスを初期化する:
  \begin{lstlisting}[language=C, numbers=none]
pthread_mutexattr_t attr;
pthread_mutexattr_init(&attr);
pthread_mutexattr_setpshared(&attr,
                             PTHREAD_PROCESS_SHARED);
pthread_mutex_init(&shm->mutex, &attr);
\end{lstlisting}
  条件変数も同様に，これらの呼び出しの \texttt{pthread\_condattr} 版で初期化
  する．\texttt{IPC\_EXCL} フラグは，セグメントがすでに存在する場合に
  \texttt{shmget} を失敗させるので，同期構成要素は一度だけ初期化される．
  その後，セグメントを接続する各プロセスは，スレッドとまったく同じように
  \texttt{shm->mutex} をロックし \texttt{shm->cond} で待つ．
\end{example}

\subsection{他の IPC 機構による同期}

プロセスは，オペレーティングシステムの IPC 機構を通じて同期することもできる．
メッセージキューは，送信と受信のプロトコルによって相互排除を提供できる
（\cref{fig:l09-mq-sync}）．P1 は共有領域にデータを書き，それから「ready」
メッセージを送る．\margin{双方向で1つのキューを使うため，各プロセスは相手の
メッセージだけを受信しなければならない．例えばメッセージの型で選択する
（System~V）．}キューでの受信でブロックしていた P2 は，そのメッセージが届くと
データを読み，「ok」メッセージで応答する．P1 は領域に再びアクセスする前に
それを待つ．各プロセスは相手のメッセージを受け取った後にのみ領域にアクセス
するので，アクセスが重なることはない．

\begin{figure}[htbp]
  \centering
  % Using a message queue to coordinate access to shared memory.
% Usage: % Using a message queue to coordinate access to shared memory.
% Usage: % Using a message queue to coordinate access to shared memory.
% Usage: \input{figures/l09-mq-sync}   (width ~116 mm)
\begin{tikzpicture}[x=1mm, y=1mm, font=\footnotesize\sffamily,
  hd/.style={draw, rounded corners=2pt, fill=ln-box, align=center,
             minimum height=9mm, inner sep=1.5pt,
             text height=1.6ex, text depth=0.5ex},
  msg/.style={font=\scriptsize\sffamily, fill=white, inner sep=0.6pt, above=0.5mm},
  lab/.style={font=\scriptsize\sffamily, text=ln-muted},
  bar/.style={draw=ln-accent, fill=ln-accent!20},
  >={Stealth[length=1.6mm]}]
  \node[hd, minimum width=16mm] (h1) at (10,0) {P1};
  \node[hd, minimum width=25mm] (hs) at (44,0) {\iflangja{共有メモリ}{shared memory}};
  \node[hd, minimum width=25mm] (hq) at (78,0) {\iflangja{メッセージキュー}{message queue}};
  \node[hd, minimum width=16mm] (h2) at (108,0) {P2};
  \foreach \x/\n in {10/h1,44/hs,78/hq,108/h2}
    \draw[dashed, ln-rule] (\n.south) -- (\x,-62);
  % blocking periods
  \draw[bar] (107,-6) rectangle (109,-28);
  \node[lab, anchor=east] at (105.5,-7) {\iflangja{\texttt{msgrcv()} で待つ}{blocked in \texttt{msgrcv()}}};
  \draw[bar] (9,-21) rectangle (11,-55);
  \node[lab, anchor=west] at (12.5,-24) {\iflangja{\texttt{msgrcv()} で待つ}{blocked in \texttt{msgrcv()}}};
  % messages
  \draw[->] (10,-12) -- node[msg] {1\ \iflangja{データを書く}{write data}} (44,-12);
  \draw[->] (10,-19) -- node[msg, pos=0.62] {2\ \texttt{msgsnd(ready)}} (78,-19);
  \draw[->] (78,-28) -- node[msg] {3\ \texttt{ready}} (108,-28);
  \draw[->] (108,-37) -- node[msg, pos=0.6] {4\ \iflangja{データを読む}{read data}} (44,-37);
  \draw[->] (108,-46) -- node[msg] {5\ \texttt{msgsnd(ok)}} (78,-46);
  \draw[->] (78,-55) -- node[msg, pos=0.45] {6\ \texttt{ok}} (10,-55);
\end{tikzpicture}
   (width ~116 mm)
\begin{tikzpicture}[x=1mm, y=1mm, font=\footnotesize\sffamily,
  hd/.style={draw, rounded corners=2pt, fill=ln-box, align=center,
             minimum height=9mm, inner sep=1.5pt,
             text height=1.6ex, text depth=0.5ex},
  msg/.style={font=\scriptsize\sffamily, fill=white, inner sep=0.6pt, above=0.5mm},
  lab/.style={font=\scriptsize\sffamily, text=ln-muted},
  bar/.style={draw=ln-accent, fill=ln-accent!20},
  >={Stealth[length=1.6mm]}]
  \node[hd, minimum width=16mm] (h1) at (10,0) {P1};
  \node[hd, minimum width=25mm] (hs) at (44,0) {\iflangja{共有メモリ}{shared memory}};
  \node[hd, minimum width=25mm] (hq) at (78,0) {\iflangja{メッセージキュー}{message queue}};
  \node[hd, minimum width=16mm] (h2) at (108,0) {P2};
  \foreach \x/\n in {10/h1,44/hs,78/hq,108/h2}
    \draw[dashed, ln-rule] (\n.south) -- (\x,-62);
  % blocking periods
  \draw[bar] (107,-6) rectangle (109,-28);
  \node[lab, anchor=east] at (105.5,-7) {\iflangja{\texttt{msgrcv()} で待つ}{blocked in \texttt{msgrcv()}}};
  \draw[bar] (9,-21) rectangle (11,-55);
  \node[lab, anchor=west] at (12.5,-24) {\iflangja{\texttt{msgrcv()} で待つ}{blocked in \texttt{msgrcv()}}};
  % messages
  \draw[->] (10,-12) -- node[msg] {1\ \iflangja{データを書く}{write data}} (44,-12);
  \draw[->] (10,-19) -- node[msg, pos=0.62] {2\ \texttt{msgsnd(ready)}} (78,-19);
  \draw[->] (78,-28) -- node[msg] {3\ \texttt{ready}} (108,-28);
  \draw[->] (108,-37) -- node[msg, pos=0.6] {4\ \iflangja{データを読む}{read data}} (44,-37);
  \draw[->] (108,-46) -- node[msg] {5\ \texttt{msgsnd(ok)}} (78,-46);
  \draw[->] (78,-55) -- node[msg, pos=0.45] {6\ \texttt{ok}} (10,-55);
\end{tikzpicture}
   (width ~116 mm)
\begin{tikzpicture}[x=1mm, y=1mm, font=\footnotesize\sffamily,
  hd/.style={draw, rounded corners=2pt, fill=ln-box, align=center,
             minimum height=9mm, inner sep=1.5pt,
             text height=1.6ex, text depth=0.5ex},
  msg/.style={font=\scriptsize\sffamily, fill=white, inner sep=0.6pt, above=0.5mm},
  lab/.style={font=\scriptsize\sffamily, text=ln-muted},
  bar/.style={draw=ln-accent, fill=ln-accent!20},
  >={Stealth[length=1.6mm]}]
  \node[hd, minimum width=16mm] (h1) at (10,0) {P1};
  \node[hd, minimum width=25mm] (hs) at (44,0) {\iflangja{共有メモリ}{shared memory}};
  \node[hd, minimum width=25mm] (hq) at (78,0) {\iflangja{メッセージキュー}{message queue}};
  \node[hd, minimum width=16mm] (h2) at (108,0) {P2};
  \foreach \x/\n in {10/h1,44/hs,78/hq,108/h2}
    \draw[dashed, ln-rule] (\n.south) -- (\x,-62);
  % blocking periods
  \draw[bar] (107,-6) rectangle (109,-28);
  \node[lab, anchor=east] at (105.5,-7) {\iflangja{\texttt{msgrcv()} で待つ}{blocked in \texttt{msgrcv()}}};
  \draw[bar] (9,-21) rectangle (11,-55);
  \node[lab, anchor=west] at (12.5,-24) {\iflangja{\texttt{msgrcv()} で待つ}{blocked in \texttt{msgrcv()}}};
  % messages
  \draw[->] (10,-12) -- node[msg] {1\ \iflangja{データを書く}{write data}} (44,-12);
  \draw[->] (10,-19) -- node[msg, pos=0.62] {2\ \texttt{msgsnd(ready)}} (78,-19);
  \draw[->] (78,-28) -- node[msg] {3\ \texttt{ready}} (108,-28);
  \draw[->] (108,-37) -- node[msg, pos=0.6] {4\ \iflangja{データを読む}{read data}} (44,-37);
  \draw[->] (108,-46) -- node[msg] {5\ \texttt{msgsnd(ok)}} (78,-46);
  \draw[->] (78,-55) -- node[msg, pos=0.45] {6\ \texttt{ok}} (10,-55);
\end{tikzpicture}

  \caption{メッセージキューが共有メモリへのアクセスを協調させる．P2 は
    「ready」の後にのみ読み，P1 は「ok」の後にのみ先へ進む．}
  \label{fig:l09-mq-sync}
\end{figure}

セマフォ（第10章）は，整数値を保持するオペレーティングシステムの同期構成要素
である．値 0 と 1 で使えばミューテックスのように振る舞う．値が 0 であるのを
見たプロセスはブロックし，値が 1 であるのを見たプロセスはそれを 0 に減らして
先へ進み，クリティカルセクションを出るときに再び値を増やす．

\begin{keypoint}
  共有メモリには，マルチスレッドプロセスの共有状態と同じ同期が必要である．
  Pthreads のミューテックスと条件変数は，プロセス共有に設定され共有領域に
  置かれていればプロセスをまたいで働く．メッセージキューとセマフォは，同じ
  協調をオペレーティングシステムを通じて提供する．
\end{keypoint}

\section{IPC のコマンドラインツール}
\label{sec:l09-tools}

System~V の IPC オブジェクトは，それを作成したプロセスより長く残るので，
Unix 系システムはそれらを調べて取り除くコマンドを提供する．\source{P3L3，IPC の
コマンドラインツール；Kerrisk ch.~45．}\texttt{ipcs} コマンドは，システム内の
すべての IPC 機構---メッセージキュー，共有メモリセグメント，セマフォ集合---を，
そのキー，識別子，所有者，アクセス権とともに一覧する．\texttt{ipcs -m} は一覧を
共有メモリに限定する．\texttt{ipcrm} コマンドは機構を取り除く．
\texttt{ipcrm -m} にセグメント識別子を続けるとそのセグメントを破棄でき，
これはプログラムが破棄し損ねたセグメントを回収する通常の方法である．

\section{共有メモリ上の通信の設計}
\label{sec:l09-design}

オペレーティングシステムは共有メモリを提供し，その後は介在しない．
\source{P3L3，共有メモリの設計上の考慮，セグメントはいくつか，設計上の考慮．}
同期機構の選択，データの形式，データを渡すプロトコルはアプリケーションに
委ねられる．この柔軟さのために，プログラマは特に2つの設計上の決定に責任を
負う．

\subsection{セグメントはいくつにするか}

共有メモリで通信するプロセスは，それを2通りに構成できる
（\cref{fig:l09-segments}）．
\begin{description}
  \item[1つの大きなセグメント] すべての通信が単一のセグメントを使う．その場合
    アプリケーションは，セグメントの一部を個々のやり取りに割り当て，後で解放
    するメモリマネージャを必要とする．
  \item[多数の小さなセグメント] 通信するプロセスの組ごとに専用のセグメントを
    使う．セグメントの作成は高価で個数にも限りがあるので，セグメントのプールを
    あらかじめ作り，空いているものの識別子をキューに保持する．プロセスはプール
    からセグメントを取り，その識別子をメッセージキューなど別の機構で相手に
    伝える．やり取りが終わるとセグメントはプールに戻る．
\end{description}

\begin{figure}[htbp]
  \centering
  % Two ways of organizing shared memory among several processes.
% Usage: % Two ways of organizing shared memory among several processes.
% Usage: % Two ways of organizing shared memory among several processes.
% Usage: \input{figures/l09-segments}   (width ~118 mm)
\begin{tikzpicture}[x=1mm, y=1mm, font=\footnotesize\sffamily,
  blk/.style={draw, minimum height=8mm, inner sep=0pt, font=\scriptsize\sffamily},
  used/.style={blk, fill=ln-accent!20},
  freeb/.style={blk, fill=white},
  box/.style={draw, rounded corners=2pt, fill=ln-box, align=center,
              inner sep=1.5pt, font=\scriptsize\sffamily},
  cell/.style={draw, fill=white, minimum width=5mm, minimum height=5mm,
               inner sep=0pt, font=\scriptsize\sffamily},
  lab/.style={font=\scriptsize\sffamily, text=ln-muted, align=center},
  >={Stealth[length=1.6mm]}]
  % (a) one large segment managed by an allocator
  \node[used,  minimum width=12mm, anchor=west] (a1) at (0,0) {P1/P2};
  \node[freeb, minimum width=8mm,  anchor=west] (a2) at (a1.east) {};
  \node[used,  minimum width=12mm, anchor=west] (a3) at (a2.east) {P1/P3};
  \node[used,  minimum width=12mm, anchor=west] (a4) at (a3.east) {P2/P3};
  \node[freeb, minimum width=8mm,  anchor=west] (a5) at (a4.east) {};
  \node[box, text width=24mm] (alloc) at (26,18)
    {\iflangja{アロケータ\\（確保・解放）}{allocator\\(allocate / free)}};
  \foreach \t in {a1,a3,a4} \draw[->] (alloc.south) -- (\t.north);
  \node[lab] at (26,-9) {\iflangja{(a) 1つの大きなセグメント}{(a) one large segment}};
  % (b) pool of small segments
  \foreach \i/\s in {1/used,2/freeb,3/used,4/freeb} {
    \node[\s, minimum width=10mm] (s\i) at (58+13*\i,0) {seg \i};
  }
  \node[lab, anchor=east] at (86.5,18) {\iflangja{空き ID}{free ids}};
  \node[cell] (q2) at (90,18) {2};
  \node[cell] (q4) at (95,18) {4};
  \draw[->, dashed] (q2.south) -- (s2.north);
  \draw[->, dashed] (q4.south) -- (s4.north);
  \node[lab] at (90.5,-9) {\iflangja{(b) 小さなセグメントのプール}{(b) pool of small segments}};
\end{tikzpicture}
   (width ~118 mm)
\begin{tikzpicture}[x=1mm, y=1mm, font=\footnotesize\sffamily,
  blk/.style={draw, minimum height=8mm, inner sep=0pt, font=\scriptsize\sffamily},
  used/.style={blk, fill=ln-accent!20},
  freeb/.style={blk, fill=white},
  box/.style={draw, rounded corners=2pt, fill=ln-box, align=center,
              inner sep=1.5pt, font=\scriptsize\sffamily},
  cell/.style={draw, fill=white, minimum width=5mm, minimum height=5mm,
               inner sep=0pt, font=\scriptsize\sffamily},
  lab/.style={font=\scriptsize\sffamily, text=ln-muted, align=center},
  >={Stealth[length=1.6mm]}]
  % (a) one large segment managed by an allocator
  \node[used,  minimum width=12mm, anchor=west] (a1) at (0,0) {P1/P2};
  \node[freeb, minimum width=8mm,  anchor=west] (a2) at (a1.east) {};
  \node[used,  minimum width=12mm, anchor=west] (a3) at (a2.east) {P1/P3};
  \node[used,  minimum width=12mm, anchor=west] (a4) at (a3.east) {P2/P3};
  \node[freeb, minimum width=8mm,  anchor=west] (a5) at (a4.east) {};
  \node[box, text width=24mm] (alloc) at (26,18)
    {\iflangja{アロケータ\\（確保・解放）}{allocator\\(allocate / free)}};
  \foreach \t in {a1,a3,a4} \draw[->] (alloc.south) -- (\t.north);
  \node[lab] at (26,-9) {\iflangja{(a) 1つの大きなセグメント}{(a) one large segment}};
  % (b) pool of small segments
  \foreach \i/\s in {1/used,2/freeb,3/used,4/freeb} {
    \node[\s, minimum width=10mm] (s\i) at (58+13*\i,0) {seg \i};
  }
  \node[lab, anchor=east] at (86.5,18) {\iflangja{空き ID}{free ids}};
  \node[cell] (q2) at (90,18) {2};
  \node[cell] (q4) at (95,18) {4};
  \draw[->, dashed] (q2.south) -- (s2.north);
  \draw[->, dashed] (q4.south) -- (s4.north);
  \node[lab] at (90.5,-9) {\iflangja{(b) 小さなセグメントのプール}{(b) pool of small segments}};
\end{tikzpicture}
   (width ~118 mm)
\begin{tikzpicture}[x=1mm, y=1mm, font=\footnotesize\sffamily,
  blk/.style={draw, minimum height=8mm, inner sep=0pt, font=\scriptsize\sffamily},
  used/.style={blk, fill=ln-accent!20},
  freeb/.style={blk, fill=white},
  box/.style={draw, rounded corners=2pt, fill=ln-box, align=center,
              inner sep=1.5pt, font=\scriptsize\sffamily},
  cell/.style={draw, fill=white, minimum width=5mm, minimum height=5mm,
               inner sep=0pt, font=\scriptsize\sffamily},
  lab/.style={font=\scriptsize\sffamily, text=ln-muted, align=center},
  >={Stealth[length=1.6mm]}]
  % (a) one large segment managed by an allocator
  \node[used,  minimum width=12mm, anchor=west] (a1) at (0,0) {P1/P2};
  \node[freeb, minimum width=8mm,  anchor=west] (a2) at (a1.east) {};
  \node[used,  minimum width=12mm, anchor=west] (a3) at (a2.east) {P1/P3};
  \node[used,  minimum width=12mm, anchor=west] (a4) at (a3.east) {P2/P3};
  \node[freeb, minimum width=8mm,  anchor=west] (a5) at (a4.east) {};
  \node[box, text width=24mm] (alloc) at (26,18)
    {\iflangja{アロケータ\\（確保・解放）}{allocator\\(allocate / free)}};
  \foreach \t in {a1,a3,a4} \draw[->] (alloc.south) -- (\t.north);
  \node[lab] at (26,-9) {\iflangja{(a) 1つの大きなセグメント}{(a) one large segment}};
  % (b) pool of small segments
  \foreach \i/\s in {1/used,2/freeb,3/used,4/freeb} {
    \node[\s, minimum width=10mm] (s\i) at (58+13*\i,0) {seg \i};
  }
  \node[lab, anchor=east] at (86.5,18) {\iflangja{空き ID}{free ids}};
  \node[cell] (q2) at (90,18) {2};
  \node[cell] (q4) at (95,18) {4};
  \draw[->, dashed] (q2.south) -- (s2.north);
  \draw[->, dashed] (q4.south) -- (s4.north);
  \node[lab] at (90.5,-9) {\iflangja{(b) 小さなセグメントのプール}{(b) pool of small segments}};
\end{tikzpicture}

  \caption{複数のプロセスの組に対する共有メモリの構成:
    (a) アロケータが一部を配る1つの大きなセグメント，
    (b) 空き識別子をキューに保持する小さなセグメントのプール．}
  \label{fig:l09-segments}
\end{figure}

\subsection{セグメントをどれだけ大きくするか}

第2の決定はセグメントの大きさと，1つに収まらないデータをどう扱うかである．
\begin{description}
  \item[データの大きさに等しいセグメント] どのデータ項目も1つのセグメントに
    収まる．データの大きさが既知で固定の場合にうまく働くが，データの大きさは
    セグメントの最大サイズに制限される．
  \item[データより小さいセグメント] より大きなデータは，一度に一部ずつ，複数の
    ラウンドに分けて転送する．その場合プロセスは，転送の進行を追うプロトコルを
    必要とする．典型的には，セグメントの先頭に置いたデータ構造が同期構成要素を
    保持し，セグメントがどれだけのデータを含むか，転送が完了したかを記録する．
\end{description}

\begin{example}
  \label{ex:l09-rounds}
  セグメントは次のヘッダで始まり，残りの部分が \SI{2}{\mega\byte} の
  ペイロードを保持する．
  \begin{lstlisting}[language=C, numbers=none]
typedef struct {
  pthread_mutex_t mutex;  /* プロセス共有 */
  pthread_cond_t  cond;   /* プロセス共有 */
  size_t nbytes;          /* buf 内のバイト数; 0 = 空 */
  int    done;            /* 最後の部分で設定 */
  char   buf[];           /* ペイロード */
} xfer_t;
\end{lstlisting}
  送信側はミューテックスを保持したまま \texttt{nbytes} が 0 になるまで待ち，
  次の部分を \texttt{buf} へコピーし，\texttt{nbytes} を設定して（最後の部分
  では \texttt{done} も設定して）signal する．受信側は \texttt{nbytes} が
  0 でなくなるまで待ち，その部分をコピーして取り出し，\texttt{nbytes} を 0 に
  戻して signal する．\SI{5}{\mega\byte} の転送は $\lceil 5/2 \rceil = 3$
  ラウンドを要する:
  \begin{center}
    \small
    \begin{tabular}{@{}cllcc@{}}
      \toprule
      ステップ & プロセス & 動作 & \texttt{nbytes}（MB） & \texttt{done} \\
      \midrule
      1 & 送信側 & バイト 0--2\,MB を書く & 2 & 0 \\
      2 & 受信側 & それを読む            & 0 & 0 \\
      3 & 送信側 & バイト 2--4\,MB を書く & 2 & 0 \\
      4 & 受信側 & それを読む            & 0 & 0 \\
      5 & 送信側 & バイト 4--5\,MB を書く & 1 & 1 \\
      6 & 受信側 & それを読み，\texttt{done} を見る & 0 & 1 \\
      \bottomrule
    \end{tabular}
  \end{center}
\end{example}

\begin{keypoint}[章のまとめ]
  IPC 機構はアドレス空間をまたいでデータを移す．メッセージパッシング---パイプ，
  メッセージキュー，ソケット---は，ポートを介してすべてのメッセージをカーネルの
  チャネルに通し，その代価としてメッセージごとにシステムコールとコピーを払う．
  共有メモリは同じ物理ページを複数のアドレス空間に対応付けるので，1回限りの
  設定の後，プロセスはカーネルを介さずに通信する．コピーとマップのどちらが安い
  かは，データの大きさと対応付けを何回再利用するかで決まり
  （\cref{eq:l09-copy-map}），Windows の LPC はこのトレードオフを利用して
  いる．System~V のセグメント（\texttt{shmget}，\texttt{shmat}，
  \texttt{shmdt}，\texttt{shmctl}）と POSIX 共有メモリ（\texttt{shm\_open}，
  \texttt{mmap}，\texttt{munmap}，\texttt{shm\_unlink}）は明示的に破棄される
  まで存続し，\texttt{ipcs} と \texttt{ipcrm} が System~V のオブジェクトを
  調べて取り除く．メモリを共有するプロセスは，領域に置いたプロセス共有の Pthreads
  ミューテックスと条件変数か，メッセージキューとセマフォによって同期しなければ
  ならず，セグメントをいくつ使うか，セグメントより大きなデータをどう転送するか
  を決めなければならない．
\end{keypoint}

\end{document}
